% Related work.
%
% Organisation follows Hu et al. (CKA-RL, NeurIPS 2025), which groups continual
% RL into the four families of the Pan et al. survey and gives each a run-in
% bold paragraph. Two families are added that the taxonomy does not cover but
% that this paper has to engage with: plasticity-oriented methods, whose
% failure mode is distinct from forgetting, and constrained RL, which supplies
% the machinery. A final paragraph places this work.
%
% Every paragraph ends on what its family leaves unaddressed, so the gap
% argument accumulates rather than arriving all at once at the end.

\section{Related Work}
\label{sec:related}

Continual reinforcement learning asks an agent to learn a sequence of tasks
while retaining competence on those it has already seen. Existing methods are
commonly grouped into four families \citep{pan2025survey}, namely
regularisation-based, rehearsal-based, architecture-based and meta-learning
based approaches. We review each in turn, then two lines of work that the
taxonomy does not cover but that bear directly on our contribution.

\paragraph{Regularisation-based methods.}
These add a penalty to the learning objective that discourages movement in
parameters deemed important to earlier tasks. Elastic weight consolidation
\citep{kirkpatrick2017ewc} weights that penalty by a Fisher-information
estimate of per-parameter importance, and Progress \& Compress
\citep{schwarz2018progress} alternates an active column with distillation into
a knowledge base under an online variant of the same penalty. More recent work
adapts the penalty during training rather than fixing it, with ITER
\citep{igl2021iter} periodically distilling into a freshly initialised network
and TRAC \citep{muppidi2024trac} tuning regularisation strength online to limit
weight drift. \citet{anand2023permanent} instead decompose the value function
into permanent and transient components to address the stability-plasticity
trade-off directly.
\textbf{What is left unaddressed.} The trade-off is expressed as a penalty
weight, and that weight has to be chosen before it is known which tasks will
come under pressure or how much. A single scalar is asked to price retention
uniformly across tasks that are not equally at risk, and the resulting
behaviour is a soft trade rather than a requirement the optimiser must meet.
Degradation on earlier tasks accordingly remains substantial as sequences
lengthen \citep{hu2025ckarl}.

\paragraph{Rehearsal-based methods.}
These retain past experience and replay it while learning new tasks. CLEAR
\citep{rolnick2019clear} combines off-policy replay with behavioural and value
cloning towards the replayed behaviour and remains the reference point for
sequential Atari evaluation \citep{powers2022cora}. RECALL
\citep{kessler2022recall} pairs multi-head training with adaptive normalisation
and policy distillation, IQ \citep{kessler2023iq} partitions contexts by online
clustering to reduce interference between state distributions, and DRAGO
\citep{zheng2024drago} substitutes synthetic replay and exploration-based
memory recovery so that raw past-task data need not be stored.
\textbf{What is left unaddressed.} The buffer is a proxy for the objective
rather than the objective itself. Replaying a fixed sample encourages the
policy to reproduce past behaviour, which is not the same as requiring that it
still achieve past return, and the two come apart as the stored distribution
drifts from what the current policy would visit. Storage also grows with the
task sequence unless the buffer is subsampled, at which point the proxy weakens
further.

\paragraph{Architecture-based methods.}
These allocate, partition or compose parameters so that tasks interfere less by
construction. Progressive networks \citep{rusu2016progressive} add a column per
task with lateral connections to frozen predecessors, PackNet
\citep{mallya2018packnet} prunes and freezes a subnetwork per task, and MaskNet
\citep{benavides2022masknet} learns per-task modulation masks over a fixed
network that can be linearly recombined for new tasks. CompoNet
\citep{malagon2024componet} makes each new policy a module that attends over
frozen previous ones, and CKA-RL \citep{hu2025ckarl} maintains a pool of
task-specific knowledge vectors, adapting a weighted combination of them to
each new task and merging similar vectors to bound growth.
\textbf{What is left unaddressed.} Protection is bought with parameters. Frozen
capacity cannot be overwritten, but neither can it be improved, so positive
backward transfer is unavailable by construction, and scalability becomes the
binding concern as the task count grows \citep{hu2025ckarl}. A related point is
evaluative rather than algorithmic. Methods in this family are normally
evaluated with a separate output head per task, which removes the most direct
channel through which a later task can overwrite an earlier one.

\paragraph{Meta-learning based methods.}
These treat the sequence as a distribution over tasks and optimise for rapid
adaptation to the next draw. MAML \citep{finn2017maml} trains parameters from
which a few gradient steps suffice on a new task, MB-MPO
\citep{clavera2018mbmpo} meta-learns policies robust to an ensemble of learned
dynamics models, and ESCP \citep{luo2022escp} strengthens context encoding so
that abrupt environment changes are absorbed quickly.
\textbf{What is left unaddressed.} The objective is adaptation speed, and
retention of specific earlier tasks is at best incidental to it. These methods
are complementary to ours rather than competing with it.

\paragraph{Plasticity-oriented methods.}
A separate failure mode is the loss of the ability to learn anything new after
prolonged training. \citet{abbas2023plasticity} show that value-based agents
cycling through Atari games progressively lose this ability as activations
sparsify and gradients shrink, and mitigate it by replacing the activation with
a concatenated ReLU. \citet{dohare2024plasticity} demonstrate the same
degradation across architectures and optimisers and propose continual
backpropagation, which continually reinitialises a small fraction of
low-utility units.
\textbf{What is left unaddressed.} These methods restore the capacity to
acquire, not the capacity to hold, and the two failures are separable. We
report this distinction empirically in Section~\ref{sec:results}, where both
interventions improve forward transfer and leave backward transfer
indistinguishable from naive fine-tuning. We make no plasticity claim, and
neither intervention conflicts with the objective proposed here.

\paragraph{Constrained RL and primal-dual optimisation.}
Constrained Markov decision processes \citep{altman1999cmdp} pose policy
optimisation subject to explicit constraints, and are most often solved either
by trust-region updates that guarantee feasibility \citep{achiam2017cpo} or by
Lagrangian relaxation with dual ascent on the multiplier
\citep{tessler2019rcpo}. The dominant application is safety, where the
constraint bounds an auxiliary cost.
\textbf{What is left unaddressed.} In that literature the constraint is an
exogenous budget, fixed in advance and independent of learning. We take the
same machinery and point it at a comparator that is learned and that moves,
namely a specialist policy trained on each task as it arrives. To our
knowledge, retention in continual RL has not previously been posed as a
constraint of this form.

\paragraph{Evaluation protocols.}
Continual RL benchmarks differ in ways that strongly affect how much
interference is possible at all. CORA \citep{powers2022cora} standardises
sequential Atari with continual evaluation, isolated forgetting and zero-shot
transfer metrics; Continual World \citep{wolczyk2021continualworld} does the
same for robotic manipulation; and the benchmark of \citet{hu2025ckarl} draws
tasks as modes of a single game and assigns each its own output head.
\textbf{What is left unaddressed.} Both a per-task head and tasks drawn from a
single game reduce the pressure on shared parameters, the first by removing the
output layer from contention and the second by making transfer plentiful. A
method that constrains retention has little room to demonstrate itself under
either. We return to this in Section~\ref{sec:setup}, where we evaluate under
settings that relax each property in turn.

\paragraph{Where this work sits.}
Across the four families, retention is approximated rather than required. It is
priced by a penalty weight, approximated by a stored sample, or purchased with
frozen parameters. This paper states it directly. We require that the deployed
policy remain within a tolerance of a task-specific specialist on every task it
has seen, and enforce that requirement by primal-dual alternation, so the
multiplier on each task is set by the optimisation rather than chosen in
advance. The constraint is one-sided, which distinguishes it from both
distillation and parameter freezing, since the deployed policy is prevented
from falling behind a specialist but left free to exceed it. Nothing is stored
and nothing is frozen, so retention does not compete with capacity as the
sequence lengthens.
